% \documentclass[a4paper, conference, compsoc]{IEEEtran}
\documentclass[12pt,a4paper]{article}
% \documentclass[%
%     pdftex,
%     oneside,			% Einseitiger Druck.
%     12pt,				% Schriftgroesse
%     parskip=half,		% Halbe Zeile Abstand zwischen Absätzen.
% %	topmargin = 10pt,	% Abstand Seitenrand (Std:1in) zu Kopfzeile [laut log: unused]
%     headheight = 12pt,	% Höhe der Kopfzeile
% %	headsep = 30pt,	% Abstand zwischen Kopfzeile und Text Body  [laut log: unused]
%     headsepline,		% Linie nach Kopfzeile.
%     footsepline,		% Linie vor Fusszeile.
%     footheight = 16pt,	% Höhe der Fusszeile
%     abstracton,		% Abstract Überschriften
%     DIV=calc,		% Satzspiegel berechnen
%     BCOR=8mm,		% Bindekorrektur links: 8mm
%     headinclude=false,	% Kopfzeile nicht in den Satzspiegel einbeziehen
%     footinclude=false,	% Fußzeile nicht in den Satzspiegel einbeziehen
%     listof=totoc,		% Abbildungs-/ Tabellenverzeichnis im Inhaltsverzeichnis darstellen
%     toc=bibliography,	% Literaturverzeichnis im Inhaltsverzeichnis darstellen
% ]{scrreprt}	% Koma-Script report-Klasse, fuer laengere Bachelorarbeiten alternativ auch: scrbook


% Basic packages
\usepackage[T1]{fontenc}
\usepackage[utf8]{inputenc}
\usepackage[scaled]{beramono}
\usepackage[english,ngerman]{babel,varioref}
\usepackage{xcolor}
\usepackage{amsmath}

% Tables
\usepackage{booktabs}
\usepackage{multirow}
\usepackage{longtable}
% Graphics and Includes
\usepackage[pdftex]{graphicx}
\graphicspath{{assets/img/}}
\DeclareGraphicsExtensions{.pdf,.jpeg,.png,.jpg}
\usepackage{tikz}
\usetikzlibrary{arrows.meta,bending,automata,shapes}
\usepackage[underline=true,rounded corners=false]{pgf-umlsd}

% Reduce distance before caption
\usepackage[skip=5pt]{caption}
% Bibliography
\usepackage[backend=biber, isbn=false, doi=false, style=ieee]{biblatex}
\addbibresource{bibliography.bib}
\AtBeginBibliography{\raggedright}
%\nocite{*}

% Gossar z.a. für acronyme
% \usepackage[acronym]{glossaries}
% \makeglossaries
\usepackage[printonlyused]{acronym} % falls gewünscht kann die Option footnote eingefügt werden, dann wird die Erklärung nicht inline sondern in einer Fußnote dargestellt

% \section*{Abkürzungsverzeichnis}
\begin{acronym}[AAAAAAA]
    \acro{wysiwyg}[WYSIWYG]{What you see is what you get}
\end{acronym}

% Colors
\definecolor{ListingBackground}{HTML}{E6E6E6}
\definecolor{LinkColor}{HTML}{000000}


% Font
\usepackage[onehalfspacing]{setspace}
\usepackage{lmodern}
\usepackage[official]{eurosym}
\usepackage{enumitem}
\usepackage[locale=DE]{siunitx} % SI Units für Währungen
\DeclareSIUnit{\EUR}{\text{\euro}} % Beispielverwendung: \SI{10.10}{\EUR}

\usepackage[autostyle=true,german=quotes]{csquotes}
\usepackage{url}
\newcommand{\code}[1]{\texttt{#1}}

% Additional Setup
\usepackage[unicode=true,hypertexnames=false,colorlinks=true,linkcolor=LinkColor,citecolor=LinkColor,urlcolor=LinkColor,pdftex]{hyperref}

% Trennung von URLs im Literaturverzeichnis (große Werte [> 10000] verhindern die Trennung)
\defcounter{biburlnumpenalty}{10} % Strafe für Trennung in URL nach Zahl
\defcounter{biburlucpenalty}{500}  % Strafe für Trennung in URL nach Großbuchstaben
\defcounter{biburllcpenalty}{500}  % Strafe für Trennung in URL nach Kleinbuchstaben
\interfootnotelinepenalty=10000 % prevent all footnotes from breaking over a page.

% Configs
\setcounter{tocdepth}{1} % Limit table of contents to subsection
\sisetup{detect-weight=true, detect-family=true} % SI Units shall detect font weight and family
\setlist[description]{style=nextline} % Break definitions of terms to a new line (used by \begin{description} \item[foo] bar \end{description})
\renewcommand*{\bibfont}{\small}

%Hurenkinder und Schusterjungen vermeiden
\clubpenalty = 10000
\widowpenalty = 10000
\displaywidowpenalty = 10000

% Quellcode
\usepackage{listings}
\usepackage{float}
\usepackage{textcomp}
\lstset{
    inputpath=assets/listings,
    language=Java,			% Standardsprache des Quellcodes
    numbers=left,			% Zeilennummern links
    stepnumber=1,			% Jede Zeile nummerieren.
    numbersep=5pt,			% 5pt Abstand zum Quellcode
    numberstyle=\tiny,		% Zeichengrösse 'tiny' für die Nummern.
    breaklines=true,		% Zeilen umbrechen wenn notwendig.
    breakautoindent=true,	% Nach dem Zeilenumbruch Zeile einrücken.
    postbreak=\space,		% Bei Leerzeichen umbrechen.
    tabsize=2,				% Tabulatorgrösse 2
    basicstyle=\ttfamily\footnotesize, % Nichtproportionale Schrift, klein für den Quellcode
    showspaces=false,		% Leerzeichen nicht anzeigen.
    showstringspaces=false,	% Leerzeichen auch in Strings ('') nicht anzeigen.
    extendedchars=true,		% Alle Zeichen vom Latin1 Zeichensatz anzeigen.
    captionpos=b,			% sets the caption-position to bottom
    backgroundcolor=\color{ListingBackground}, % Hintergrundfarbe des Quellcodes setzen.
    xleftmargin=0pt,		% Rand links
    xrightmargin=0pt,		% Rand rechts
    frame=single,			% Rahmen an
    frameround=ffff,
    rulecolor=\color{darkgray},	% Rahmenfarbe
    fillcolor=\color{ListingBackground},
    keywordstyle=\color[rgb]{0.133,0.133,0.6},
    commentstyle=\color[rgb]{0.133,0.545,0.133},
    stringstyle=\color[rgb]{0.627,0.126,0.941},
    float,
    upquote=true
}


\colorlet{punct}{red!60!black}
\definecolor{delim}{RGB}{20,105,176}
\colorlet{numb}{magenta!60!black}

\lstdefinelanguage{json}{
    string=[s]{"}{"},
    stringstyle=\color{numb},
    literate=
     *{0}{{{\color{numb}0}}}{1}
      {1}{{{\color{numb}1}}}{1}
      {2}{{{\color{numb}2}}}{1}
      {3}{{{\color{numb}3}}}{1}
      {4}{{{\color{numb}4}}}{1}
      {5}{{{\color{numb}5}}}{1}
      {6}{{{\color{numb}6}}}{1}
      {7}{{{\color{numb}7}}}{1}
      {8}{{{\color{numb}8}}}{1}
      {9}{{{\color{numb}9}}}{1}
      {:}{{{\color{punct}{:}}}}{1}
      {,}{{{\color{punct}{,}}}}{1}
      {\{}{{{\color{delim}{\{}}}}{1}
      {\}}{{{\color{delim}{\}}}}}{1}
      {[}{{{\color{delim}{[}}}}{1}
      {]}{{{\color{delim}{]}}}}{1}
}

\lstdefinelanguage{yaml}{
  identifierstyle=\color{delim},
  sensitive=false,
  comment=[l]{\#},
  commentstyle=\color{purple}\ttfamily,
  stringstyle=\color{numb},
  morestring=[b]',
  morestring=[b]"
}

\lstloadlanguages{Python,Java,bash}



% Algorithmen (Pseudocode)
\usepackage[german,onelanguage,linesnumbered,ruled]{algorithm2e}
\newcommand\mycommfont[1]{\color{gray}\ttfamily{#1}}
\SetCommentSty{mycommfont}
% Useful tools
\usepackage{blindtext}
\usepackage{lipsum}
\usepackage[section]{placeins} % Prevent figures and tables to float to new section
\usepackage[obeyFinal,backgroundcolor=yellow,linecolor=black]{todonotes}


%Angaben zur Arbeit
\def\myTitel{Erlernen von Kundenpräferenzen in einer Wohnungsbörse}
\def\myArbeit{Programmentwurf}
\def\myFach{Maschinelles Lernen}
\def\myDatum{\today}
\def\myBetreuer{Prof. Dr. Dirk Reichardt}
\def\myGutachter{Strenger Gutachter}
\def\myBearbeitungszeit{Lange}
\def\myAbgabeort{Heilbronn}

%Angaben zur Person
\def\myAutor{Daniel Knecht, Lea Poletin}
\def\myDegree{Master Informatik}
\def\myInstitution{DHBW CAS, Bildungscampus 13, 74076 Heilbronn}
\def\myMatrikelnr{123456, 8514262}
\def\myKurs{TMINF17}
\def\myFirma{SAP SE}
\def\myFirmenort{Walldorf}
%\section*{Abkürzungsverzeichnis}
\begin{acronym}[AAAAAAA]
    \acro{wysiwyg}[WYSIWYG]{What you see is what you get}
\end{acronym}

\begin{document}


\title{\myTitel}
\author{\myAutor~(\myMatrikelnr)\\
    \small{\myFach}\\
    \small{\myDegree}\\
    \small{\myInstitution}
}
\date{\myDatum}

\begin{titlepage}
	\begin{longtable}{p{8.2cm} p{5.4cm}}
		% Firmenlogo
		&
		\includegraphics[width=5.4cm]{casLogo}
	\end{longtable}
	\addtocounter{table}{-1}
    {\let\newpage\relax\maketitle}
    \thispagestyle{empty}
    % \vfill
    \begin{abstract}
Dieser Programmentwurf befasst sich mit der Problemstellung einer Wohnungsbörse.
Es werden Konzepte aus dem Bereich des \emph{Maschinellen Lernen} entwickelt
und eines dieser Konzepte implementiert.
Die Systeme erlernen die gegebene Bewertungsfunktion
umn anschließend neue Datensätze nach dem gleichen Schema zu bewerten.
\end{abstract}
\end{titlepage}
% \begin{titlepage}
	\begin{longtable}{p{8.2cm} p{5.4cm}}
		% Firmenlogo
		&
		\includegraphics[width=5.4cm]{casLogo}
	\end{longtable}
	\addtocounter{table}{-1}
	\enlargethispage{20mm}
	\begin{center}
		\vspace*{12mm}	{\LARGE\textbf \myTitel }\\
		\vspace*{12mm}	{\large\textbf \myArbeit}\\
		% \vspace*{12mm}	\langdeckblattabschlusshinleitung\\
		\vspace*{3mm}		{\textbf \myDegree}\\
		\vspace*{12mm}	des Studiengangs \myKurs\\
    \vspace*{3mm}		am DHBW CAS\\
		\vspace*{12mm}	von\\
		\vspace*{3mm}		{\large\textbf \myAutor}\\
		\vspace*{12mm}	\myDatum\\
	\end{center}
	\vfill
	\begin{spacing}{1.2}
	\begin{tabbing}
		mmmmmmmmmmmmmmmmmmmmmmmmmm             \= \kill
		\textbf{Bearbeitungszeit}       \>  \myBearbeitungszeit\\
		\textbf{Matrikelnummer, Kurs}  \>  \myMatrikelnr, \myKurs\\
		\textbf{Firma}                  \>  \myFirma, \myFirmenort\\
		\textbf{Betreuer}               \>  \myBetreuer\\
		\textbf{Gutachter}              \>  \myGutachter
	\end{tabbing}
	\end{spacing}
\end{titlepage}

\pagenumbering{Roman}

% Sperrvermerk
\thispagestyle{empty}
% Sperrvermerk direkt hinter Titelseite
\section*{Sperrvermerk}

\vspace*{2em}

Die vorliegende {\myArbeit} mit dem Titel {\itshape{} \myTitel{}\/}
enthält unternehmensinterne bzw. vertrauliche Informationen der {\myFirma},
ist deshalb mit einem Sperrvermerk versehen
und wird ausschließlich zu Prüfungszwecken am Studiengang
{\myDegree} der Dualen Hochschule Baden-Württemberg vorgelegt.
Sie ist ausschließlich zur Einsicht durch den zugeteilten Gutachter,
die Leitung des Studiengangs und ggf. den Prüfungsausschuss des Studiengangs
bestimmt.  Es ist untersagt,
\begin{itemize}
\item den Inhalt dieser Arbeit (einschließlich Daten, Abbildungen, Tabellen, Zeichnungen usw.) als Ganzes oder auszugsweise weiterzugeben,
\item Kopien oder Abschriften dieser Arbeit (einschließlich Daten, Abbildungen, Tabellen, Zeichnungen usw.) als Ganzes oder in Auszügen anzufertigen,
\item diese Arbeit zu veröffentlichen bzw. digital, elektronisch oder virtuell zur Verfügung zu stellen.
\end{itemize}
Jede anderweitige Einsichtnahme und Veröffentlichung – auch von Teilen der Arbeit – bedarf der vorherigen Zustimmung durch den Verfasser und {\myFirma}.

\vspace{3em}

\myAbgabeort, \myDatum
\vspace{4em}

\rule{6cm}{0.4pt}\\
\myAutor

\newpage

% Erklärung
\thispagestyle{empty}

\section*{Selbstständigkeitserklärung}
\vspace*{2em}


Wir versicheren hiermit,
dass der {\myArbeit} mit dem Thema: {\itshape \myTitel } selbstständigverfasst und keine anderen
als die angegebenen Quellen und Hilfsmittel verwendet wurden.
Wir versicheren zudem, dass die eingereichte elektronische Fassung mit der gedruckten Fassung übereinstimmt.


\vspace{3em}
\noindent
\myAbgabeort, \myDatum
\\
\\
\noindent
\begin{tabular}{l l}
\rule{6cm}{0.4pt} &  \rule{6cm}{0.4pt}\\
\myAutorA         &  \myAutorB

\end{tabular}

\newpage


% \begin{abstract}
Dieser Programmentwurf befasst sich mit der Problemstellung einer Wohnungsbörse.
Es werden Konzepte aus dem Bereich des \emph{Maschinellen Lernen} entwickelt
und eines dieser Konzepte implementiert.
Die Systeme erlernen die gegebene Bewertungsfunktion
umn anschließend neue Datensätze nach dem gleichen Schema zu bewerten.
\end{abstract}
% \newpage

\pagestyle{plain}		% nur Seitenzahlen im Fuß

% Inhaltsverzeichnis
\begin{spacing}{1.1}
    \begingroup
        \pagestyle{empty}
        %für die Anzeige von Unterkapiteln im Inhaltsverzeichnis
        \setcounter{tocdepth}{2}

        \tableofcontents
        \clearpage
    \endgroup
\end{spacing}
\newpage

% Abkürzungsverzeichnis
\cleardoublepage
\section*{Abkürzungsverzeichnis}
\begin{acronym}[AAAAAAA]
    \acro{wysiwyg}[WYSIWYG]{What you see is what you get}
\end{acronym}


% Abbildungsverzeichnis
\cleardoublepage
\listoffigures

%Tabellenverzeichnis
\cleardoublepage
\listoftables

% Quellcodeverzeichnis
\cleardoublepage
\lstlistoflistings
\cleardoublepage

\pagenumbering{arabic}

\section{Einleitung}\label{sec:einleitung}
In diesem Programmentwurf wird das Problem \emph{Erlernen von Kundenpräferenzen in einer Wohnungsbörse}
bearbeitet. Gegeben sind zwei Datensätze: 
\begin{itemize}
    \item Datensatz zum Trainieren mit erwartetem Ergebnis
    \item Datensatz zum Testen ohne erwartetes Ergebnis
\end{itemize}

Die Datensätze beschreiben 22 Merkmale von Wohnungen die in einer Wohnungsbörse angeboten werden. \\
Zunächst werden die gegebenen Daten analysiert. Anschließend werden die aus der Analyse gewonnenen Informationen 
in ein Konzept für einen Entscheidungsbaum einfließen. 
Als ein alternatives Konzept wird der Entwurf eines neuronalen Netzes vorgestellt. Hierfür werden unteranderem
die Daten vorverarbeitet und die Struktur und Arbeitsweise des neuronalen Netzes definiert. Das Netz wird allerdings nicht praktisch getestet. 
Um das Problem der Wohnungsbörse zu lösen wird abschließend der Entscheidungsbaum implementiert und optimiert. 
Der Entscheidungsbaum wird mit den Trainingsdaten aufgebaut und mit den Testdaten validiert.
Die Ergebnisse und Erkenntnisse werden schließlich zusammengefasst und bewertet.
\section{Kapitel}\label{sec:kaptiel}

\blindtext

\begin{figure}[h]
    \centering
    \includegraphics[width=0.5\textwidth]{placeholder}
    \caption{a nice placeholder}
    \label{fig:place}
\end{figure}

\blindtext
Außerdem lässt sich so eine \autoref{fig:place} referenzieren\cite{dreyfus:1980}.


\lstinputlisting[
    float,
    floatplacement=H,
    caption=Code-Beispiel aus \cite{mustermann:2012},
    label=lst:hello
]{hello-world.java}

Auch lässt sich das \autoref{lst:hello} referenzieren.
Wobei folgendes zu beachten ist:
\blindtext

\lstinputlisting[
    float,
    floatplacement=H,
    caption=Python-Code,
    label=Python-Code
    language=Python
]{quicksort.py}
\input{content/03_Fazit}

%\printacronyms{}
\printbibliography{}
\end{document}
